\begin{filecontents*}[overwrite]{mixed_mahler_verify.py}
import json
import math
import sympy as sp

z,X,Y,C,W=sp.symbols('z X Y C W')
P=X**3-X+z
E=(1+z)*Y**2-1
tauP=P.subs({z:z**3,X:X-z}, simultaneous=True)
tauE=E.subs({z:z**3,Y:Y/(1+z)}, simultaneous=True)
assert sp.cancel(tauP-(P-3*z*X*(X-z)))==0
assert sp.cancel(tauE-(E-3*z/(1+z)**2*(E+1)))==0
assert sp.expand((-z*X**2+z**2*X).subs(z,X-X**3)-(X**7-X**5))==0

def v3fact(n):
    s=0
    while n:
        n//=3;s+=n
    return s
def weight(n): return n+v3fact(n)
def w9(n):
    s=0
    while n:
        s+=n;n//=9
    return s
table=[]
for g in (1,2,3,4,10,100):
    d=next(n for n in range(1,g+1) if weight(n)>=g)
    count=6*sum(weight(i)+weight(j)<g for i in range(d) for j in range(d))
    e=next(n for n in range(1,g+1) if w9(n)>=g)
    table.append([g,d,6*d*d,count,6*g,6*e])
assert [r[3] for r in table]==[6,18,36,48,216,14646]

# Independent coefficient arithmetic, always truncated modulo z**N.
N=512; modulus=3**12
zero=[0]*N
one=zero.copy();one[0]=1
zz=zero.copy();zz[1]=1
def add(*terms): return [sum(t[i] for t in terms)%modulus for i in range(N)]
def scale(a,k): return [x*k%modulus for x in a]
def mul(a,b):
    out=[0]*N
    aa=[(i,x) for i,x in enumerate(a) if x]
    bb=[(i,x) for i,x in enumerate(b) if x]
    for i,x in aa:
        for j,y in bb:
            if i+j>=N: break
            out[i+j]=(out[i+j]+x*y)%modulus
    return out
def power(a,n):
    out=one
    while n:
        if n&1:out=mul(out,a)
        n//=2
        if n:a=mul(a,a)
    return out
def invplus(k):return [((-1)**n*math.comb(k+n-1,n))%modulus for n in range(N)]
A=zero.copy();q=1
while q<N:A[q]=1;q*=3
B=one.copy();q=1
while q<N:
    for i in range(N-1,q-1,-1):B[i]=(B[i]+B[i-q])%modulus
    q*=3
PA=add(power(A,3),scale(A,-1),zz)
EB=add(mul(add(one,zz),power(B,2)),scale(one,-1))
assert all(x%3==0 for x in PA+EB)
CC=[x//3 for x in PA];WW=[x//3 for x in EB]
U=add(power(PA,3),scale(PA,-9),scale(mul(zz,mul(A,add(A,scale(zz,-1)))),27))
V=add(power(EB,3),scale(EB,-9),scale(mul(mul(zz,invplus(2)),add(EB,one)),27))
assert all(x%81==0 for x in U+V)
Cf=add(power(CC,3),scale(CC,-1),mul(zz,power(A,2)),scale(mul(power(zz,2),A),-1))
Wf=add(power(WW,3),scale(WW,-1),mul(zz,invplus(2)))
assert all(x%3==0 for x in Cf+Wf)
H=add(A,B)
MH=add(power(add(power(H,3),scale(H,-1),zz),2),scale(mul(power(zz,2),invplus(3)),-1))
assert all(x%3==0 for x in MH)
Z=[x//3 for x in MH]
rhs=add(mul(mul(mul(zz,invplus(1)),B),CC),mul(mul(zz,invplus(2)),WW),mul(mul(mul(zz,invplus(2)),A),H))
assert all((x-y)%3==0 for x,y in zip(Z,rhs))

# Resultant checks at convergent 3-adic evaluation points.
resultant_checks=[]
modulus2=3**11
for z0 in (3,6,9,12):
    aa=0;bb=1;q=1
    while q<64:
        t=pow(z0,q,modulus2);aa=(aa+t)%modulus2;bb=bb*(1+t)%modulus2;q*=3
    hh=(aa+bb)%modulus2
    inv1=pow(1+z0,-1,modulus2);inv3=pow(1+z0**3,-1,modulus2)
    p1=(X**3-X+z0)**3-9*((X-z0)**3-(X-z0)+z0**3)
    q1=((hh-X)**2-inv1)**3-9*((hh-X)**2*inv1**2-inv3)
    p1=sp.Poly(p1,X,modulus=modulus2).as_expr()
    q1=sp.Poly(q1,X,modulus=modulus2).as_expr()
    residual=int(sp.resultant(p1,q1,X))%modulus2
    assert residual%(3**10)==0
    resultant_checks.append([z0,residual])

print(json.dumps({'symbolic_identities':'passed','coefficients_checked':N,'series_checks':'passed','table':table,'resultant_mod_3_power_11':resultant_checks},indent=2))
\end{filecontents*}
\documentclass[11pt,a4paper]{article}
\usepackage[T1]{fontenc}
\usepackage[utf8]{inputenc}
\usepackage{lmodern,microtype}
\usepackage[margin=27mm]{geometry}
\usepackage{amsmath,amssymb,amsthm,mathtools}
\usepackage{booktabs,array,enumitem,xcolor}
\usepackage[colorlinks=true,linkcolor=blue!45!black,citecolor=blue!45!black,urlcolor=blue!45!black]{hyperref}
\usepackage{fancyhdr}
\pagestyle{fancy}\fancyhf{}
\fancyhead[L]{\small Mixed Mahler congruences}
\fancyhead[R]{\small KEIO}
\fancyfoot[C]{\thepage}
\setlength{\headheight}{14pt}
\setlength{\parskip}{3pt}
\setlength{\parindent}{1.2em}
\setlist{itemsep=2pt,topsep=4pt}
\newtheorem{theorem}{Theorem}[section]
\newtheorem{lemma}[theorem]{Lemma}
\newtheorem{proposition}[theorem]{Proposition}
\newtheorem{corollary}[theorem]{Corollary}
\theoremstyle{definition}
\newtheorem{definition}[theorem]{Definition}
\newtheorem{question}[theorem]{Question}
\theoremstyle{remark}
\newtheorem{remark}[theorem]{Remark}
\newcommand{\Z}{\mathbb Z}
\newcommand{\Q}{\mathbb Q}
\newcommand{\C}{\mathbb C}
\newcommand{\F}{\mathbb F}
\newcommand{\A}{\Phi_3}
\newcommand{\B}{\Psi}
\newcommand{\val}{\nu_3}
\DeclareMathOperator{\Res}{Res}
\DeclareMathOperator{\Gal}{Gal}
\DeclareMathOperator{\trdeg}{trdeg}
\DeclareMathOperator{\ord}{ord}
\DeclareMathOperator{\Frac}{Frac}
\hypersetup{pdftitle={Mixed Mahler Congruences: Exact Low-Precision Kernels and Optimal Relations for a Lacunary Sum plus a Cantor Product},pdfauthor={KEIO}}
\title{Mixed Mahler Congruences\\[3pt]
\large Exact low-precision kernels, improved norm bounds,\\
and optimal relations for a lacunary sum plus a Cantor product}
\author{KEIO\thanks{Research draft prepared with AI assistance. The proofs and finite checks are supplied for independent review; the manuscript is not externally peer reviewed or proof-assistant certified.}}
\date{26 October 2026}
\begin{document}
\maketitle
\begin{abstract}
Let $A(z)=\sum_{n\ge0}z^{3^n}$ and
$B(z)=\prod_{n\ge0}(1+z^{3^n})$. We study their joint algebra modulo
powers of $3$ over $R=\Z[z,z^{-1},(1+z)^{-1}]$.
The individual minimum-degree problems of Krattenthaler and M\"uller
provide a starting point, but do not determine mixed relations.
We give exact joint evaluation kernels modulo $3,9,27,81$, with
respectively $6,18,36,48$ necessary and sufficient module generators.
For $H=A+B$, the minimum degree of a monic annihilator is exactly
$6\gamma$ for $1\le\gamma\le9$, and exactly $54$ for $\gamma=10$.
The lower bounds use two independent Artin--Schreier correction classes.
For every polynomial $F(A,B)$, a root-counting norm construction gives
an annihilator of degree at most $6d$, where
$d+\sum_{j\ge1}\lfloor d/9^j\rfloor\ge\gamma$.
This is $(16/3)\gamma+O(\log\gamma)$, improving the elementary $6\gamma$
bound. We also give an all-precision mixed spanning theorem, classify
the reduced mixed fields for every odd prime, and supply a direct
characteristic-zero independence proof. The full mixed kernel and the
optimal degree for $H$ beyond precision $3^{10}$ are left open.
\end{abstract}

\section{Introduction and precise scope}
\label{sec:intro}

Two complementary series used in prime-power congruence methods are
\begin{equation}\label{eq:AB}
 A(z)=\sum_{n\ge0}z^{3^n},\qquad
 B(z)=\prod_{n\ge0}(1+z^{3^n}).
\end{equation}
They satisfy the exact identities
\begin{equation}\label{eq:functional}
 A(z^3)=A(z)-z,\qquad B(z^3)=\frac{B(z)}{1+z}.
\end{equation}
Throughout, congruence of Laurent series means coefficientwise
congruence. Our fixed coefficient ring is
\[
 R=\Z[z,z^{-1},(1+z)^{-1}]\subset\Z((z)).
\]
For $H\in\Z((z))$ admitting monic relations modulo every $3^\gamma$, write
\begin{equation}\label{eq:mu}
 \mu_H(\gamma)=\min\{\deg_T G:G\in R[T]\text{ monic},\
                         G(H)\equiv0\pmod{3^\gamma}\}.
\end{equation}
For a Laurent series or an element of $R$, $\val$ is the least $3$-adic
valuation of its coefficients; $\val(0)=\infty$.
The notation $R[A,B]\bmod3^\gamma$ always denotes the image under
coefficientwise reduction, equivalently $R[X,Y]/J_\gamma$ for the
evaluation ideal defined below. It does not denote the intrinsic
quotient $R[A,B]/3^\gamma R[A,B]$; coefficientwise divisibility need
not be divisibility inside that subalgebra.

Krattenthaler and M\"uller posed two separate minimum-degree conjectures.
Their Conjecture~4 in \cite{KMp} concerns $\Phi_p(z)=\sum z^{p^n}$ in
$\Z[z,T]$: the proposed minimum is the least $D$ with $\nu_p(D!)\ge\gamma$.
Their Conjecture~12 in \cite{KM3} concerns $B$ in the ring $R$: the proposed
minimum is $2d_\gamma$, where
\begin{equation}\label{eq:w}
 w(d)=d+\val(d!),\qquad d_\gamma=\min\{d\ge1:w(d)\ge\gamma\}.
\end{equation}
The arithmetic part of \cite{KEIO70} supplies a Frobenius construction
attaining both prior lower bounds. Section~\ref{sec:lifting} reproduces
the construction and its coefficient-ring descent so that the present
mixed results do not depend on an unexamined assertion in that draft.

The substantive joint assertions in this paper are the following.
\begin{enumerate}
\item The evaluation ideal in $R[X,Y]$ is computed exactly through
precision $81$ (Theorem~\ref{thm:kernel}). The associated counts are
minimal numbers of generators, not free ranks.
\item For $H=A+B$,
\begin{equation}\label{eq:mainexact}
 \mu_H(\gamma)=
 \begin{cases}6\gamma,&1\le\gamma\le9,\\54,&\gamma=10.
 \end{cases}
\end{equation}
See Theorem~\ref{thm:sumexact}. This does not follow by multiplying the
two individual minimum degrees.
\item Set
\begin{equation}\label{eq:W9}
 W_9(d)=\sum_{j\ge0}\left\lfloor\frac d{9^j}\right\rfloor,
 \qquad e_\gamma=\min\{d\ge1:W_9(d)\ge\gamma\}.
\end{equation}
For every $F\in R[X,Y]$, the series $F(A,B)$ has a monic annihilator
of degree $6e_\gamma$ modulo $3^\gamma$
(Theorem~\ref{thm:norm}). No affine functional equation for $F(A,B)$
is assumed.
\end{enumerate}

We separate these joint conclusions from standard consequences of
Galois theory and Mahler's method, which are included to explain the
structure. We do not identify or claim the resolution of an additional
previously published conjecture by the mixed results. In particular,
the Ap\'ery congruence conjectures at the end of \cite{KM3} were already
settled in \cite{Apéry}. The questions in Section~\ref{sec:open} are
questions posed here, not relabelled historical conjectures.

\section{Individual Frobenius constructions and descent}
\label{sec:lifting}

\subsection{The lifting identity}
Let $p$ be a prime and let a coefficient ring of integral Laurent series
be stable under $\sigma:c(z)\mapsto c(z^p)$. Suppose
$f(z^p)=a(z)f(z)+b(z)$, and define
\[
 \tau G(z,T)=G(z^p,a(z)T+b(z)).
\]

\begin{lemma}[Frobenius correction]\label{lem:lift}
If $G$ is monic of degree $D>0$ and $G(z,f(z))$ is divisible by $p^e$,
then
\[
 G^p-p^{(p-1)e}\tau G
\]
is monic of degree $pD$ and its evaluation is divisible by $p^{pe+1}$.
Starting with degree $D_0$ and divisibility $p^{e_0}$ gives degrees
$D_0p^r$ and divisibilities
$p^{e_r}$, where $e_r=p^re_0+(p^r-1)/(p-1)$.
\end{lemma}
\begin{proof}
Write $G(f)=p^e C(z)$. Evaluation commutes with the substitutions, so
\[
 (G^p-p^{(p-1)e}\tau G)(f)
 =p^{pe}\bigl(C(z)^p-C(z^p)\bigr).
\]
The parenthesis is divisible by $p$ by the coefficientwise Frobenius
congruence. The correction has degree at most $D<pD$, which preserves
monicity. The recurrence $e_{r+1}=pe_r+1$ gives the formula.
\end{proof}

If $d=\sum d_rp^r$ is a base-$p$ expansion, multiplying the corresponding
tower polynomials to powers $d_r$ gives degree $D_0d$ and valuation at least
\begin{equation}\label{eq:digitval}
 \sum_r d_re_r=e_0d+\nu_p(d!).
\end{equation}
The digit-product step was already described in the original papers;
it is the recursively available tower that supplies all the inputs.

For $\Phi_p$, take the seed $T$, $e_0=0$, and $a=1,b=-z$.
This gives monic polynomials of every degree $i$, with evaluation
valuation at least $\nu_p(i!)$. Combined with the lower bound in
\cite[Lemma~1]{KMp}, this attains the value in its Conjecture~4.
The construction applies to every prime, not only $3$.

\subsection{The Cantor ring and its finite descent}
For $B$, initially work in
\[
 R_\infty=\Z[z,z^{-1},(1+z^{3^j})^{-1}:j\ge0].
\]
This ring is stable under $z\mapsto z^3$. Frobenius and
\eqref{eq:functional} imply $B^2\equiv(1+z)^{-1}\pmod3$.
Start with $T^2-(1+z)^{-1}$, degree $2$ and valuation at least $1$.
Lemma~\ref{lem:lift} and \eqref{eq:digitval} give a monic polynomial
of degree $2d$ with evaluation valuation at least $w(d)$.

We must return to the prescribed ring $R$, rather than silently enlarging
it. Put $U_j=1+z^{3^j}$, $V_j=(1+z)^{3^j}$, and
$D_j=(U_j-V_j)/3\in\Z[z]$. For every $N\ge1$,
\begin{equation}\label{eq:descent}
 U_j^{-1}\equiv
 \sum_{a=0}^{N-1}(-3D_j)^a V_j^{-a-1}\pmod{3^N}.
\end{equation}
Indeed the exact remainder is
$(-3D_j)^N V_j^{-N}U_j^{-1}$. Substitution in each of the finitely
many coefficients descends a polynomial from $R_\infty$ to $R$ to any
fixed precision. Its leading coefficient $1$ is retained unchanged.
The lower bound of \cite[Theorem~10]{KM3} now gives the asserted minimum
$2d_\gamma$ in its Conjecture~12.

\begin{proposition}[Two triangular families]\label{prop:bases}
There are fixed monic polynomials $a_i\in R[X]$ and $b_j\in R[Y]$,
of degrees $i,j$, such that
\[
 \val(a_i(A))\ge\alpha_i:=\val(i!),\qquad
 \val(b_j(B))\ge\beta_j:=w(\lfloor j/2\rfloor).
\]
They may be chosen algorithmically by the preceding towers and finite
descent, using precision $\beta_j+1$ for each $b_j$.
\end{proposition}
\begin{proof}
The first construction was just given. For $j=2d$, use the descended
Cantor polynomial; for $j=2d+1$, multiply it by $Y$.
Since $B$ is a unit, the valuation lower bound is unchanged.
\end{proof}

Only these lower bounds, not an unproved tensor-product independence
assertion, are used in the next section. Also
\begin{equation}\label{eq:alphaidentity}
 \alpha_{3r+s}=w(r)\quad(0\le s<3),\qquad
 \beta_{2j+\epsilon}=w(j)\quad(0\le\epsilon<2).
\end{equation}

\section{Mixed spanning at arbitrary precision}

\begin{theorem}[Weighted mixed spanning]\label{thm:spanning}
Modulo $3^\gamma$, every polynomial in $A,B$ with coefficients in $R$
is an $R$-linear combination of
\[
 a_i(A)b_j(B),\qquad \alpha_i+\beta_j<\gamma.
\]
Consequently it is generated by at most
\begin{equation}\label{eq:Ngamma}
 N_\gamma=6\#\{(r,j)\in\Z_{\ge0}^2:w(r)+w(j)<\gamma\}
\end{equation}
elements. As $\gamma\to\infty$,
\[
 N_\gamma=\frac43\gamma^2+O(\gamma\log\gamma).
\]
\end{theorem}
\begin{proof}
Monicity makes the products $a_i(X)b_j(Y)$ an $R$-basis of $R[X,Y]$
by triangular change of basis from the monomials. Terms whose weight
is at least $\gamma$ vanish upon evaluation modulo $3^\gamma$.
Equation~\eqref{eq:alphaidentity} gives the count.
Legendre's formula gives $w(r)=(3r-s_3(r))/2=(3/2)r+O(\log(r+2))$.
The permitted pairs therefore differ by only an $O(\log\gamma)$-wide
boundary strip from the integer points in $r+j<2\gamma/3$.
The strip contains $O(\gamma\log\gamma)$ points; the triangle contains
$2\gamma^2/9+O(\gamma)$ points. Multiplication by $6$ proves the formula.
\end{proof}

Separate monic reductions instead give a rectangle of
$(3d_\gamma)(2d_\gamma)=6d_\gamma^2$ terms, asymptotic to
$(8/3)\gamma^2$. Some sample counts are:
\begin{center}
\begin{tabular}{rrrrr}
\toprule
$\gamma$ & $d_\gamma$ & Rectangle & $N_\gamma$ & $6\gamma$\\
\midrule
1&1&6&6&6\\
2&2&24&18&12\\
3&3&54&36&18\\
4&3&54&48&24\\
10&8&384&216&60\\
100&69&28566&14646&600\\
\bottomrule
\end{tabular}
\end{center}
The last column is an elementary \emph{annihilator-degree} bound for
one mixed element, not a count of simultaneous module generators.
These are different optimization problems. In fact, a monic relation
of degree $6$ exists modulo $3$ for every mixed element, and its
$\gamma$-th power gives the last column. Section~\ref{sec:norm} improves
that bound. We assert sharpness of $N_\gamma$ only for $\gamma\le4$.

\section{The reduced mixed field}
\label{sec:reduced}

The following description also clarifies why separate reductions do not
settle the mixed problem. We state it for every odd prime $p$.
Let $x,y$ be the coefficientwise reductions of
$\sum z^{p^n}$ and $\prod(1+z^{p^n})$, and put $K=\F_p(z)$.
They satisfy
\begin{equation}\label{eq:charp}
 x^p-x=-z,\qquad y^{p-1}=(1+z)^{-1}.
\end{equation}

\begin{theorem}[Mixed degree classification]\label{thm:classification}
The field $L=K(x,y)$ is Galois over $K$, of degree $p(p-1)$, with
automorphisms $x\mapsto x+c$, $y\mapsto\zeta y$ for
$c\in\F_p$, $\zeta\in\F_p^\times$.
For a nonzero $f\in L$, write uniquely
\[
 f=\sum_{j=0}^{p-2} f_j(x)y^j,\qquad f_j\in K[X],\quad\deg f_j<p.
\]
Set $d=\gcd(p-1,\{j:f_j\ne0\})$, $m=(p-1)/d$, and let
$\epsilon=0$ if every $f_j$ is constant in $X$, and $\epsilon=1$ otherwise.
Then
\[
 [K(f):K]=p^\epsilon m,\qquad
 K(f)=\begin{cases}K(y^d),&\epsilon=0,\\K(x,y^d),&\epsilon=1.
 \end{cases}
\]
The genus of this field is $0$ if $\epsilon=0$, and
$(p-1)(m-1)/2$ if $\epsilon=1$.
\end{theorem}
\begin{proof}
The Artin--Schreier polynomial for $x$ is irreducible: a rational function
$u^p-u$ cannot have a pole of order $1$ at infinity. Here the absence
of a rational root suffices: after adjoining one root all the roots
are obtained by adding elements of $\F_p$, so its splitting-field
Galois group is a subgroup of the prime-order group $(\F_p,+)$.
The polynomial
for $y^{-1}$ is Eisenstein at $z+1$, of degree $p-1$. The two Galois
degrees are coprime, so their compositum has the asserted group.

Translation invariance of a polynomial in $x$ of degree less than $p$
forces it to be constant, as its highest nonconstant term has nonzero
finite difference. The Kummer subgroup fixing $f$ has exactly $d$
elements. Every subgroup of the direct product of these coprime cyclic
groups splits into its two components. Thus the stabilizer has order
$p^{1-\epsilon}d$, proving the degree and the stated fixed field.

If $\epsilon=0$, then $z=(y^d)^{-m}-1$, so the field is rational.
Otherwise set $U=x$ and $V=y^{-d}$; its equation is
$V^m=1+U-U^p$. The right side has $p$ simple roots and a pole of order
$p$ at infinity. The tame Kummer cover is fully ramified at these
$p+1$ points. Riemann--Hurwitz gives
$2g-2=-2m+(p+1)(m-1)$, as asserted.
\end{proof}

\begin{corollary}[Sum and product]\label{cor:minmodp}
The following are irreducible minimal polynomials over $K$:
\[
 \mu_{x+y}(T)=(T^p-T+z)^{p-1}-\frac{z^{p-1}}{(1+z)^p},
\]
\[
 \mu_{xy}(T)=(T^p-qT)^{p-1}-z^{p-1}q^p,
 \qquad q=(1+z)^{-1}.
\]
Both have degree $p(p-1)$.
\end{corollary}
\begin{proof}
For the sum, $(x+y)^p-(x+y)+z=-zy/(1+z)$.
For the product, $(xy)^p=qxy-zqy$.
Taking $(p-1)$-st powers gives the displayed relations.
The degree classification, or the distinct conjugates, proves minimality.
\end{proof}

For $p=3$, put $h=x+y$ and retain $K=\F_3(z)$, $L=K(x,y)$.
Then $[L:K]=6$, $L=K(h)$, and
\begin{equation}\label{eq:M}
 M(T)=(T^3-T+z)^2-\frac{z^2}{(1+z)^3}\in R[T]
\end{equation}
reduces to the minimal polynomial of $h$. This degree-$6$ field has
genus $1$; its proper intermediate fields have genus $0$.

\section{Independent first corrections and exact mixed kernels}
\label{sec:kernel}

Define
\begin{equation}\label{eq:PE}
 P(X)=X^3-X+z,\qquad E(Y)=(1+z)Y^2-1,
\end{equation}
and define integral power series and their reductions by
\[
 C_0=\frac{P(A)}3,\quad W_0=\frac{E(B)}3,
 \qquad c=\overline{C_0},\quad w=\overline{W_0}.
\]
Lowercase $w$ here is a series; the weight function $w(d)$ always has an
integer argument.

\begin{lemma}[Two independent correction classes]\label{lem:AS}
Over $L$, the elements $c,w$ satisfy
\begin{equation}\label{eq:AS}
 c^3-c=-zx^2+z^2x=x^7-x^5=:f,
 \qquad w^3-w=-\frac{z}{(1+z)^2}=:g.
\end{equation}
Their compositum has degree $9$ over $L$, with Galois group $\F_3^2$
acting by independent translations of $c,w$.
\end{lemma}
\begin{proof}
Direct substitution in \eqref{eq:functional} gives
\[
 C_0(z^3)=C_0(z)-zA(A-z),
\]
\[
 W_0(z^3)=W_0(z)-\frac{z}{(1+z)^2}
                  -\frac{3z}{(1+z)^2}W_0(z).
\]
Reduction modulo $3$ gives \eqref{eq:AS}. Since $z=x-x^3$, its first
right side equals $x^7-x^5$.

It remains to check independence of the two Artin--Schreier classes.
We may extend constants to an algebraic closure for this check.
The field $L$ has equation $y^{-2}=1+x-x^3$. At the place over
$x=\infty$, $x$ has pole order $2$, hence $f$ has pole order $14$;
$g$ is regular there. At any place over $1+z=0$, the quadratic Kummer
extension has ramification index $2$, so $g$ has pole order $4$;
$f$ is regular there. Both $14$ and $4$ are prime to $3$.
A function of the form $u^3-u$ has pole order divisible by $3$ at every
pole. Thus no nonzero $\F_3$-linear combination of $f,g$ is of that
form: use infinity if the coefficient of $f$ is nonzero, and a place
over $-1$ otherwise. The two classes are independent. Equivalently,
their two cyclic degree-$3$ extensions are distinct and linearly
disjoint, giving the asserted degree and group.
\end{proof}

\begin{lemma}[Exact valuation in a finite box]\label{lem:box}
For arbitrary $d_{s,e,r,j}\in R$, set
\[
 D(X,Y)=\sum_{\substack{0\le s,r,j\le2\\0\le e\le1}}
 d_{s,e,r,j}X^sY^eP(X)^rE(Y)^j.
\]
Then
\begin{equation}\label{eq:boxval}
 \val(D(A,B))=
 \min_{s,e,r,j}\bigl(\val(d_{s,e,r,j})+r+j\bigr).
\end{equation}
In particular the formula is exact for every polynomial with
$\deg_X<9$, $\deg_Y<6$, written in this triangular basis.
\end{lemma}
\begin{proof}
The zero polynomial is immediate. Otherwise divide the evaluation by
the power of $3$ given by the minimum, and
reduce modulo $3$. The terms attaining the minimum become a linear
combination of $x^sy^ec^rw^j$. The six elements $x^sy^e$ form a
$K$-basis of $L$, and the nine elements $c^rw^j$ form an $L$-basis
by Lemma~\ref{lem:AS}. Their $54$ products are $K$-independent,
so the reduced combination is nonzero. Division of elements of $R$
by their exact powers of $3$ stays in $R$, since its inverted
denominators have nonzero reductions. This proves the equality.
\end{proof}

Put $\mathfrak m=(3,P,E)\subset R[X,Y]$ and
\begin{align}
 U(X)&=P(X)^3-9P(X)+27zX(X-z),\label{eq:UV1}\\
 V(Y)&=E(Y)^3-9E(Y)
       +\frac{27z}{(1+z)^2}\bigl(E(Y)+1\bigr).\label{eq:UV2}
\end{align}
Their evaluations are divisible by $81$, by \eqref{eq:AS} or the
exact functional identities in its proof. The polynomial $U$ is
monic of degree $9$ in $X$; $V$ has degree $6$ in $Y$ and leading
coefficient $(1+z)^3$, a unit in $R$.

\begin{theorem}[Exact kernels through $81$]\label{thm:kernel}
Let
\[
 J_\gamma=\{G\in R[X,Y]:G(A,B)\equiv0\pmod{3^\gamma}\}.
\]
Then
\begin{equation}\label{eq:kernels}
 J_1=\mathfrak m,\qquad J_2=\mathfrak m^2,\qquad
 J_3=\mathfrak m^3,\qquad J_4=\mathfrak m^4+(U,V).
\end{equation}
For $1\le\gamma\le4$, the evaluated algebra has the following
decomposition as an $R$-module:
\begin{equation}\label{eq:module}
 R[A,B]\bmod3^\gamma\ \cong\
 \bigoplus_{\substack{0\le r,j\le2\\r+j<\gamma}}
       \bigl(R/3^{\gamma-r-j}R\bigr)^6.
\end{equation}
Its minimum number of $R$-module generators is respectively
$6,18,36,48$.
\end{theorem}
\begin{proof}
The displayed generators of the ideals have the required evaluation
divisibility, which proves one inclusion. Divide an arbitrary polynomial
by $U$ and by $(1+z)^{-3}V$ to obtain a remainder of degrees $<9,<6$.
These divisions commute since each divisor uses only its own polynomial
variable. Both discarded terms evaluate to multiples of $81$.
Expand the remainder as in Lemma~\ref{lem:box}. Its evaluation vanishes
modulo $3^\gamma$ precisely when each coefficient is divisible by
$3^{\max(\gamma-r-j,0)}$. Every resulting term belongs to
$\mathfrak m^\gamma$. Also $U,V\in\mathfrak m^3$, so for
$\gamma\le3$ the discarded terms already belong to
$\mathfrak m^\gamma$; for $\gamma=4$ they are the two additional
generators. This proves \eqref{eq:kernels}.

The same exact coefficient criterion proves the module decomposition.
After quotienting this module by $3$ and extending scalars from
$R/3R$ to $K$, its dimension is the number of summands. Any generating
set must have at least that many elements, and the displayed basis
classes attain the count. These are cyclic torsion summands, not a
claim that the algebra is a free module over $R/3^\gamma R$.
\end{proof}

\section{A norm construction at every precision}
\label{sec:norm}

We first prove a root-counting lemma that is useful beyond the present
pair of series. The valuation argument takes place in a complete
discretely valued field $\mathcal K$ of characteristic $0$, with
valuation ring $\mathcal O$, uniformizer $3$, and an isometric
endomorphism $\sigma$ inducing Frobenius on the residue field.
For our series one can take the fraction field of
\[
 \mathcal O=\varprojlim_n(\Z/3^n\Z)((z)),
\]
with $\sigma(z)=z^3$. Thus $A,B,R_\infty$ embed in $\mathcal O$,
and their original coefficientwise valuation is unchanged.

\begin{lemma}[Root counts along a Frobenius tower]\label{lem:packets}
Suppose $a\in\mathcal O^\times$, $b,F\in\mathcal O$, and
$\sigma F=aF+b$. Let $G_0\in\mathcal O[T]$ be monic, with
$v(G_0(F))\ge1$ and $G_0'(F)\in\mathcal O^\times$. Put
\[
 e_r=\frac{3^{r+1}-1}{2},\qquad
 G_{r+1}=G_r^3-3^{2e_r}\tau G_r,
 \quad\tau G_r(T)=(\sigma G_r)(aT+b).
\]
For every $1\le t\le r+1$, exactly $3^{r-t+1}$ roots $\alpha$
of $G_r$, counted with multiplicity, satisfy $v(\alpha-F)\ge t$.
All roots are integral.
\end{lemma}
\begin{proof}
Monicity and integrality give the last assertion. Define
\[
 v_{r,t}=\sum_{j=1}^t3^{r-j+1},\qquad
 Z_{r,t}(Y)=3^{-v_{r,t}}G_r(F+3^tY).
\]
We prove that $Z_{r,t}$ is integral and its reduction has degree
$3^{r-t+1}$, and that for $t=r+1$ this reduction is affine with
unit linear coefficient. For $r=0,t=1$, this follows by Taylor
expansion and the two hypotheses on $G_0$.

For $t\le r$, the recurrence gives
\[
 Z_{r+1,t}=Z_{r,t}^3-
       3^{2(e_r-v_{r,t})}(\sigma Z_{r,t})(aY).
\]
Here $e_r-v_{r,t}\ge1$ and $v_{r+1,t}=3v_{r,t}$, so its reduction
is the cube of the previous reduction, of the required degree.
At $t=r+1$, we have $v_{r,t}=e_r$. Write
$\overline Z_{r,r+1}=h+uY$, with $u\ne0$. Scaling the next polynomial
by $3^{3e_r}$ gives reduction
\[
 (h+uY)^3-(h^3+u^3\bar aY)
       =u^3Y^3-u^3\bar aY.
\]
It has degree $3$ and unit linear coefficient. Substitute $Y\mapsto3Y$
and divide by $3$ for the next depth. Its constant coefficient is
integral: putting $C=3^{-e_r}G_r(F)\in\mathcal O$ gives
$G_{r+1}(F)=3^{3e_r}(C^3-\sigma C)\in3^{3e_r+1}\mathcal O$.
Its linear coefficient is a unit and
the higher coefficients vanish on reduction. This proves the induction.

For completeness, a primitive integral polynomial has exactly as many
integral roots as the degree of its reduction, counted with multiplicity.
In a finite valued splitting extension, factor it as
\[
 c\prod_{v(\alpha)\ge0}(Y-\alpha)
  \prod_{v(\alpha)<0}(1-Y/\alpha),
\]
absorbing the factors $-\alpha$ into $c$. Every displayed factor is
integral and primitive. Their reductions have degrees $1$ and $0$,
respectively. The product is primitive because its reduction is
nonzero, so $c$ is a unit. Its reduced degree therefore counts exactly
the integral roots.
Apply it to $Z_{r,t}$. Integral values of
$Y=(\alpha-F)/3^t$ are exactly the roots in the asserted ball.
\end{proof}

For the pair $A,B$, take the two seed polynomials
\[
 G_0(X)=X^3-X+z,\qquad Q_0(Y)=Y^2-(1+z)^{-1},
\]
and the substitutions $X\mapsto X-z$, $Y\mapsto Y/(1+z)$.
Both derivatives at the corresponding series are units modulo $3$.
Use Lemma~\ref{lem:packets} to form $G_r,Q_r$; they are monic of
degrees $3\cdot3^r$ and $2\cdot3^r$ over $R_\infty$.

\begin{theorem}[Mixed norm bound]\label{thm:norm}
For every $F\in R[X,Y]$ and $r\ge0$ there is a monic polynomial
$N_{r,F}\in R[T]$ of degree $6\cdot9^r$ such that
\begin{equation}\label{eq:normval}
 \val\bigl(N_{r,F}(F(A,B))\bigr)\ge\frac{9^{r+1}-1}{8}.
\end{equation}
Consequently
\begin{equation}\label{eq:allupper}
 \mu_{F(A,B)}(\gamma)\le6e_\gamma
 =\frac{16}{3}\gamma+O(\log\gamma),
\end{equation}
where $e_\gamma$ is defined in \eqref{eq:W9}.
\end{theorem}
\begin{proof}
The quotient
\[
 \mathcal A_r=R_\infty[X,Y]/(G_r(X),Q_r(Y))
\]
is free of rank $6\cdot9^r$. Take the characteristic polynomial
$\widetilde N_{r,F}(T)$ of multiplication by $F$ on this module.
It is monic with coefficients in $R_\infty$. In a splitting field it is
\[
 \prod_{\alpha:G_r(\alpha)=0}\ \prod_{\beta:Q_r(\beta)=0}
           (T-F(\alpha,\beta)),
\]
with multiplicities. Polynomial evaluation with integral coefficients
is valuation-Lipschitz:
\[
 v\bigl(F(A,B)-F(\alpha,\beta)\bigr)
 \ge\min\{v(A-\alpha),v(B-\beta)\}.
\]
For each integer depth $t=1,\ldots,r+1$, the root-counting lemma
gives $9^{r-t+1}$ pairs with both distances at least $t$.
All other factors have nonnegative valuation. Summing these contributions
gives
\[
 v\bigl(\widetilde N_{r,F}(F(A,B))\bigr)
 \ge\sum_{t=1}^{r+1}9^{r-t+1}=\frac{9^{r+1}-1}{8}.
\]
Descend the finitely many coefficients by \eqref{eq:descent} to this
precision, preserving monicity and degree, to obtain $N_{r,F}\in R[T]$.

For $d=\sum d_r9^r$, multiply $N_{r,F}^{d_r}$.
Its degree is $6d$ and its valuation is at least
\[
 \sum_r d_r\frac{9^{r+1}-1}{8}
 =\sum_{j\ge0}\left\lfloor\frac d{9^j}\right\rfloor=W_9(d).
\]
Choose $d=e_\gamma$. Finally $W_9(d)=(9d-s_9(d))/8$
implies $e_\gamma=(8/9)\gamma+O(\log\gamma)$.
\end{proof}

\begin{remark}
This is a univariate annihilator bound for each fixed mixed element.
It improves $6\gamma$ beginning at $\gamma=10$: the first two tower
points are degree $54$ at precision $3^{10}$ and degree $486$ at
precision $3^{91}$. It is not a proof that $6e_\gamma$ is minimal
at arbitrary precision, nor a complexity bound for expanding the
resultant or determinant.
\end{remark}

\section{Optimal relations for the sum through \texorpdfstring{$3^{10}$}{3 to the tenth power}}
\label{sec:sum}

Let $H=A+B$ and let $M$ be the monic degree-$6$ polynomial
in \eqref{eq:M}. Its evaluation is divisible by $3$; put
$Z_0=M(H)/3$ and $v=\overline{Z_0}$. The next lemma detects the
first genuinely mixed obstruction.

\begin{lemma}[The sum correction is primitive]\label{lem:primitive}
Let $a=(1+z)^{-1}$ and $h=x+y$. Then
\begin{equation}\label{eq:vlinear}
 v=zay\,c+za^2w+za^2xh.
\end{equation}
Moreover $[L(v):L]=9$ and $[K(h,v):K]=54$.
\end{lemma}
\begin{proof}
In characteristic zero use $B^2=a(1+3W_0)$ and set
\[
 J=C_0+aBW_0+ABH.
\]
Then $H^3-H+z=-zaB+3J$, whence the exact identity
\[
 \frac{M(H)}3=z^2a^3W_0-2zaBJ+3J^2.
\]
On reducing modulo $3$, use $y^2=a$ and $-2=1$ to obtain
\eqref{eq:vlinear}. An automorphism $(i,j)\in\F_3^2$ of
$L(c,w)/L$ translates $v$ by $zay\,i+za^2j$.
The ratio of these two coefficients is $y/a=(1+z)y$, whose square
is $1+z$. It is not in $\F_3$. Therefore no nonzero pair $(i,j)$
fixes $v$, so its orbit has size $9$. Finally $L=K(h)$ has degree $6$.
\end{proof}

\begin{proposition}[Exact valuation below degree $54$]\label{prop:sumval}
Every $G\in R[T]$ of degree less than $54$ has a unique expansion
\[
 G(T)=\sum_{r=0}^{8} g_r(T)M(T)^r,\qquad\deg g_r<6.
\]
Writing $\val(g_r)$ for the minimum valuation of its coefficients,
\begin{equation}\label{eq:sumval}
 \val(G(H))=\min_{0\le r\le8}\{\val(g_r)+r\}.
\end{equation}
\end{proposition}
\begin{proof}
The expansion is monic polynomial division. Divide the evaluation by
the least indicated power of $3$ and reduce. The resulting combination
is in the $54$ elements $h^sv^r$, $0\le s<6$, $0\le r<9$.
They are independent over $K$ by Lemma~\ref{lem:primitive}, so at
least one surviving coefficient prevents cancellation.
\end{proof}

\begin{theorem}[The exact initial degree sequence]\label{thm:sumexact}
For $H=A+B$, formula \eqref{eq:mainexact} holds in the original ring $R$.
\end{theorem}
\begin{proof}
For $1\le\gamma\le9$, the polynomial $M^\gamma$ is a monic
annihilator of degree $6\gamma$. If a monic $G$ had degree
$D<6\gamma\le54$, its highest nonzero $M$-coefficient would be monic
and have index $\lfloor D/6\rfloor\le\gamma-1$.
Proposition~\ref{prop:sumval} would give
$\val(G(H))\le\gamma-1$, a contradiction.

For $\gamma=10$, Theorem~\ref{thm:norm} with $r=1$ provides a monic
degree-$54$ annihilator. Every monic polynomial of degree less than
$54$ has valuation at most $8$ by the same argument, so no smaller
degree is possible.
\end{proof}

An explicit compact prescription for the degree-$54$ polynomial is
\begin{align*}
 G_1(X)&=(X^3-X+z)^3
             -9\bigl((X-z)^3-(X-z)+z^3\bigr),\\
 Q_1(Y)&=\left(Y^2-\frac1{1+z}\right)^3
             -9\left(\frac{Y^2}{(1+z)^2}-\frac1{1+z^3}\right),\\
 \widetilde N(T)&=\Res_X\bigl(G_1(X),Q_1(T-X)\bigr).
\end{align*}
This is monic of degree $54$. Descent of $(1+z^3)^{-1}$ by
\eqref{eq:descent} with $N=10$, before or after taking the resultant,
places its coefficients in $R$ without changing its evaluation modulo
$3^{10}$. The construction does not require printing a large expanded
polynomial.

\begin{center}
\begin{tabular}{rrrrrrrrrrr}
\toprule
$\gamma$&1&2&3&4&5&6&7&8&9&10\\
\midrule
$\mu_H(\gamma)$&6&12&18&24&30&36&42&48&54&54\\
\bottomrule
\end{tabular}
\end{center}
The different first improvement thresholds are informative: individual
relations improve on their naive powers at $\gamma=4$, while the
mixed sum does not improve on $M^\gamma$ until $\gamma=10$.

\section{Characteristic-zero independence and specialization}
\label{sec:independence}

The arithmetic dependence modulo prime powers should not be confused
with dependence in characteristic zero. The next argument is a direct
specialization of established asymptotic ideas in Mahler theory. The
underlying independence also falls under classical mixed
additive--multiplicative criteria; see \cite{Kubota,Vaananen}.

\begin{theorem}[A direct radial proof]\label{thm:independence}
For an integer $b\ge2$ that is not a power of $2$, the functions
\[
 A_b(z)=\sum_{n\ge0}z^{b^n},\qquad
 B_b(z)=\prod_{n\ge0}(1+z^{b^n})
\]
are algebraically independent over $\C(z)$.
\end{theorem}
\begin{proof}
Set $z=e^{-t}$, $t\downarrow0$, and
$N=\lfloor\log_b(1/t)\rfloor$. Summing
$1-e^{-b^nt}=O(b^nt)$ for $n<N$ and the exponentially decaying tail
for $n\ge N$ gives
\[
 A_b(e^{-t})=\frac{\log(1/t)}{\log b}+O(1).
\]
Similarly $\log(1+e^{-b^nt})=\log2+O(b^nt)$ for $n<N$, and the
remaining logarithms have a uniformly bounded sum. Thus, with
$\theta=\log_b2$, there are positive constants $c,C$ such that
\[
 ct^{-\theta}\le B_b(e^{-t})\le Ct^{-\theta}.
\]

Suppose a nonzero polynomial relation exists, and clear denominators:
\[
 \sum_{j=0}^J q_j(z,A_b(z))B_b(z)^j=0,
       \qquad q_j\in\C[z,X].
\]
For each nonzero $q_j=\sum_i q_{ij}(z)X^i$, let
$m_j=\min_i\ord_{z=1}q_{ij}$ and choose the largest $i=d_j$
attaining that minimum. Taylor expansion at $1$ and the first
asymptotic estimate give
\[
 q_j(e^{-t},A_b(e^{-t}))
    =c_jt^{m_j}(\log(1/t))^{d_j}(1+o(1)),\quad c_j\ne0.
\]
Because $\theta$ is irrational, the exponents $m_j-j\theta$ are
pairwise distinct. The term with the smallest exponent dominates all
the others in absolute value, even allowing the bounded multiplicative
oscillation of $B_b$. Dividing the relation by that nonzero term gives
$1+o(1)=0$, a contradiction. Finally $\log_b2$ is rational for an
integer $b\ge2$ exactly when $b$ is a power of $2$.
\end{proof}

\begin{corollary}[Values at regular algebraic points]\label{cor:values}
For $b$ as above and every algebraic $\alpha$ with $0<|\alpha|<1$,
$A_b(\alpha)$ and $B_b(\alpha)$ are algebraically independent over
$\overline\Q$. In particular, for $b=3$ their sum, product, and every
nonconstant polynomial in the two values with algebraic coefficients
are transcendental.
\end{corollary}
\begin{proof}
Apply Nishioka's specialization theorem, in the form of
\cite[Theorem~1]{AF}, to
\[
 \begin{pmatrix}1\\A_b(z)\\B_b(z)\end{pmatrix}
 =\begin{pmatrix}1&0&0\\z&1&0\\0&0&1+z\end{pmatrix}
 \begin{pmatrix}1\\A_b(z^b)\\B_b(z^b)\end{pmatrix}.
\]
The matrix and its inverse are regular at every iterate
$\alpha^{b^n}$, since these points are nonzero and have modulus less
than $1$. Its determinant is $1+z$. The functional transcendence
degree is $2$ by Theorem~\ref{thm:independence}, so the same is true
for the values. A nonconstant algebraic-coefficient polynomial value
could not be algebraic without providing a nonzero polynomial relation.
\end{proof}

This corollary is an application of a prior theorem, not an additional
claimed solution of an open transcendence conjecture. It concerns
formal-series values inside the unit disk, and has no implication for
the Mills-type constants considered elsewhere in the corpus.

\section{What remains open and what has been established}
\label{sec:open}

The proof of the low-precision kernel uses two independent first
corrections. Higher corrections need further independence arguments;
neither the two individual conjectures nor characteristic-zero
independence supplies them automatically.

\begin{question}[Higher mixed evaluation kernels]
Determine $J_\gamma$ for $\gamma\ge5$. Is there a compatible choice
of the triangular families in Proposition~\ref{prop:bases} for which
every finite expansion satisfies
\[
 \val\left(\sum_{i,j} c_{ij}a_i(A)b_j(B)\right)
 =\min_{i,j}\bigl(\val(c_{ij})+\alpha_i+\beta_j\bigr)?
\]
\end{question}
The equality, if proved with the coefficient-ring details retained,
would make the all-precision spanning count $N_\gamma$ sharp. The
inequality in the direction $\ge$ is already proved here. We do not
assume the reverse inequality outside the finite box of
Lemma~\ref{lem:box}.

\begin{question}[The optimal degree sequence of the sum]
Does $\mu_{A+B}(\gamma)=6e_\gamma$ hold for every $\gamma\ge1$?
\end{question}
Theorem~\ref{thm:norm} proves the upper bound at every precision and
Theorem~\ref{thm:sumexact} proves equality for $\gamma\le10$.
For $\gamma>10$ this paper provides no matching general lower bound.
An equality assertion in this range would require new work, not an
extrapolation from the displayed table.

\begin{question}[Other primitive mixed elements]
For which $F\in R[X,Y]$ is the norm bound optimal? How does the
degree of the first normalized correction over $K(\bar F)$ determine
the initial segment of $\mu_{F(A,B)}$?
\end{question}
The sum calculation gives one complete example: the reduced field
has degree $6$ and the first correction has degree $9$ over it.
The universal norm construction explains why base $9$, rather than
base $3$, appears in the resulting mixed upper bound.

For attribution, the one-series lower bounds belong to
Krattenthaler--M\"uller, the Frobenius construction and descent were
developed in \cite{KEIO70}, and the Galois, Newton-polygon, and Mahler
specialization tools are classical. The contributions developed here
are the explicit mixed application of those tools, including the
exact kernels, the degree-$54$ obstruction and construction, and the
all-precision mixed norm bound. A scoped literature check did not
locate identical joint statements; this is not an assertion of an
exhaustive priority search or external acceptance.

\appendix
\section{Exact checks and reproducibility}
\label{app:checks}

The proofs are infinite-series and field-theoretic arguments. Finite
checks are supplementary and cannot establish a minimum-degree lower
bound for arbitrary coefficient functions.

Compiling this single TeX source also writes
\texttt{mixed\_mahler\_verify.py}, an optional verification script requiring
Python~3 and SymPy. Run it with
\texttt{python3 mixed\_mahler\_verify.py}.
The script checks the symbolic identities below, the first $512$
coefficients of the two correction equations and the sum-correction
identity, the $81$-divisibility of $U(A),V(B)$ at that truncation,
and the integer counting tables. It also checks the degree-$54$
resultant modulo $3^{10}$ at the convergent $3$-adic evaluation points
$z=3,6,9,12$. These checks passed for the present draft; the checks
at particular values do not prove the formal-series assertion.

The following identities are convenient independent checks of signs
and denominators. All substitutions are simultaneous:
\begin{align*}
 P(z^3,X-z)&=P(z,X)-3zX(X-z),\\
 E\left(z^3,\frac{Y}{1+z}\right)
  &=E(z,Y)-\frac{3z}{(1+z)^2}\bigl(E(z,Y)+1\bigr),\\
 z=x-x^3&\ \Longrightarrow\ -zx^2+z^2x=x^7-x^5.
\end{align*}
They can be verified by polynomial expansion over $\Q(z)$.
The first two also verify the displayed formulas for $U,V$ and the
Artin--Schreier equations. The exact identity preceding
\eqref{eq:vlinear} verifies the primitive-correction computation.

For the numerical tables, the following complete Python code uses only
integer arithmetic and reproduces the counting formulas. It does not
test or replace the independence proofs.
\begin{quote}\begin{minipage}{\linewidth}\small\ttfamily
\begin{tabbing}
xxxx\=xxxx\=xxxx\=\kill
def v3fact(n):\\
\>s = 0\\
\>while n:\\
\>\>n //= 3\\
\>\>s += n\\
\>return s\\
def weight(n):\\
\>return n + v3fact(n)\\
def w9(n):\\
\>s = 0\\
\>while n:\\
\>\>s += n\\
\>\>n //= 9\\
\>return s\\
for g in (1, 2, 3, 4, 10, 100):\\
\>d = next(n for n in range(1, g+1) if weight(n) >= g)\\
\>N = 6 * sum(weight(i)+weight(j) < g\\
\>\>for i in range(d) for j in range(d))\\
\>e = next(n for n in range(1, g+1) if w9(n) >= g)\\
\>print(g, d, 6*d*d, N, 6*g, 6*e)
\end{tabbing}
\end{minipage}\end{quote}
The last two columns compare the elementary and norm degree upper
bounds. At $\gamma=100$ they are $600$ and $540$; this is an upper-bound
comparison, not a computed exact minimum.

\begin{thebibliography}{99}
\bibitem{KMp}
C. Krattenthaler and T. W. M\"uller,
\emph{A method for determining the mod-$p^k$ behaviour of recursive sequences},
Aequationes Math. \textbf{99} (2025), 2505--2547.
\href{https://doi.org/10.1007/s00010-025-01195-0}{doi:10.1007/s00010-025-01195-0}.
Preprint first posted 2015; revised version:
\href{https://arxiv.org/abs/1508.02580v2}{arXiv:1508.02580v2}.

\bibitem{KM3}
C. Krattenthaler and T. W. M\"uller,
\emph{A method for determining the mod-$3^k$ behaviour of recursive sequences},
preprint (2013),
\href{https://arxiv.org/abs/1308.2856}{arXiv:1308.2856}.
Theorem~10 and Conjecture~12 refer to this version.

\bibitem{KEIO70}
KEIO,
\emph{Frobenius Lifting, Floquet Rigidity, and Equivariant Stabilization},
research manuscript [070], Verifiable Mathematical Reasoning Corpus,
24 October 2026, arithmetic sections T76--T78.
\href{https://github.com/QEEDDD/Verifiable-Mathematical-Reasoning-Corpus}{Public corpus}.
Source inspected at commit \texttt{dd9c67cb72c5991cce5816de540db8c2132ca343}.
Unreviewed research draft.

\bibitem{Apéry}
C. Krattenthaler and T. W. M\"uller,
\emph{Generalised Ap\'ery numbers modulo $9$},
J. Number Theory \textbf{147} (2015), 708--720.
\href{https://arxiv.org/abs/1401.1444}{arXiv:1401.1444}.

\bibitem{Kubota}
K. K. Kubota,
\emph{On the algebraic independence of holomorphic solutions of certain
functional equations and their values},
Math. Ann. \textbf{227} (1977), 9--50.

\bibitem{Vaananen}
K. V\"a\"an\"anen,
\emph{Algebraic independence of certain Mahler numbers},
preprint (2015),
\href{https://arxiv.org/abs/1507.02510}{arXiv:1507.02510}.
See Theorem K for the mixed additive--multiplicative criterion.

\bibitem{AF}
B. Adamczewski and C. Faverjon,
\emph{A new proof of Nishioka's theorem in Mahler's method},
C. R. Math. \textbf{361} (2023), 1011--1028.
\href{https://doi.org/10.5802/crmath.458}{doi:10.5802/crmath.458}.
\end{thebibliography}
\end{document}
